\section{What the finite-sample diagnostic can establish}
\label{app:statistical}
Theorem~\ref{thm:spectral-proof} requires population $W_2$ errors, or
justified upper bounds on them. The sample-based sliced distances and
finite-feature discrepancies do not supply those upper bounds.
We therefore ask a separate statistical question: do the numerical and
target distributions agree on a fixed collection of measurements?

Each measurement is the mean of a bounded test function, called a witness.
The finite-witness metric is the largest absolute difference between the
numerical and target means. It is a finite integral probability metric.
Concentration bounds give confidence intervals for this particular metric.
Simultaneous coverage means that the confidence statement applies to all
reported comparisons together. Such testing and kernel mean comparisons
are standard~\citep{gretton2012kernel}.\AIassistance

\begin{proposition}[Simultaneous target-fidelity intervals]
\label{prop:statistical-proof}
Condition on all training data, fitted fields, diagnostic times,
solver configurations, and a finite dictionary $f_1,\ldots,f_J$ with
$|f_j|\le1$. For every numerical or target law, estimate its witness means
from $n_l$ independent draws, independent of this conditioning. Let $Q_{\rm
means}$ be the total number of scalar means estimated. For failure
probability $0<\delta<1$, set
\[
 r_l=\sqrt{\frac{2\log(2Q_{\rm means}/\delta)}{n_l}}.
\]
With probability at least $1-\delta$, simultaneously for all reported
comparisons of law $l$ and target $p$,
\begin{align}
 L_l&=\max\{0,\widehat D_l-r_l-r_p\}\le D_{\cal F}(\mu_l,p),\\
 D_{\cal F}(\mu_l,p)&\le U_l=\min\{2,\widehat D_l+r_l+r_p\},
 \label{eq:statistical-interval}
\end{align}
where $D_{\cal F}(\mu,\nu)=\max_j|\E_\mu f_j-\E_\nu f_j|$ and
$\widehat D_l$ is its empirical counterpart. Samples can be shared between
resolutions, times, and target comparisons; independence across these
means is not required. Independence within each sample mean is required.
If each witness is also 1-Lipschitz, $L_l\le W_1(\mu_l,p)\le W_2(\mu_l,p)$.
\end{proposition}
\begin{proof}
First bound the error of each sample mean. Then add the failure
probabilities to bound the chance that any one of those bounds fails.
This is the union-bound step. It gives simultaneous coverage of every
reported comparison and does not require different comparisons to be independent.

For a variable in $[-1,1]$, Hoeffding's bounded-mean inequality gives
$\Pr(|\widehat m-m|>r)\le2\exp(-nr^2/2)$.
This follows, for example, by bounding the centered moment generating
function by $\exp(s^2/2)$, applying Markov's inequality, and minimizing in
$s$; the lower tail is identical. A union bound over all means gives
simultaneous errors at most $r_l$. Triangle inequalities for each difference
and then the finite maximum give \eqref{eq:statistical-interval}.
The final Wasserstein inequality is the Lipschitz coupling bound.
\end{proof}

Hoeffding's interval uses only the witness's range. It does not use how
much the observed values vary. An empirical Bernstein interval also uses
the sample variance and can therefore be narrower. For this alternative,
pair each numerical sample with one independent target sample and write $D_i=f_j(\widehat X_i)-f_j(Y_i)\in[-2,2]$.
Let $v_n$ be its unbiased sample variance. For $Q_{\rm diff}$ scalar
differences and $n\ge2$, a simultaneous two-sided radius is
\begin{equation}
 r^{\rm EB}=\sqrt{\frac{2v_n\log(4Q_{\rm diff}/\delta)}{n}}
     +\frac{28\log(4Q_{\rm diff}/\delta)}{3(n-1)}.
 \label{eq:bernstein-radius}
\end{equation}
This is the empirical Bernstein inequality of Maurer and
Pontil~\citep[Theorem 4]{maurer2009bernstein}, scaled from $[0,1]$ to
$[-2,2]$. We combine the failure probabilities for both signs and every
selected difference, so that the intervals hold simultaneously. For each metric,
take the maximum of $\max\{0,|\widehat D_j|-r_j^{\rm EB}\}$ as its lower
bound and the maximum of $\min\{2,|\widehat D_j|+r_j^{\rm EB}\}$ as its
upper bound. The experiment splits its total failure probability equally between
this family and the Hoeffding family so both displays hold simultaneously.

If $L_{\rm fine}>U_{\rm coarse}$, the intervals establish worsening in
this \emph{specified finite-witness metric} at the stated confidence level.
They do not establish worsening in full $W_2$, because $U_{\rm coarse}$
is not a Wasserstein upper bound.
Small finite-witness discrepancy does not establish distributional fidelity.
Two distributions can agree on every selected mean while differing on a test
that was not included in the dictionary. A low cross-grid discrepancy alone is
still weaker: the zero field has perfectly consistent outputs at every
resolution while following the wrong path whenever the target moves.
Neither test identifies an asymptotic order from finitely many resolutions.

The confidence statement assumes that the evaluation set does not select
the models or witnesses. Selecting them on that set requires either a
suitable complexity bound or a new evaluation set. A finite collection of
models specified in advance is allowed if every mean is counted in
$Q_{\rm means}$. The experimental dictionary comprises fixed sine
and cosine functions with frequency norm at most one, chosen independently
of all training and evaluation samples. The intervals can be wide at 1,024 samples. An inconclusive interval
establishes neither an increase nor its absence. The comparison on the larger evaluation set
uses the same fitted fields with independent evaluation samples.

\section{Global regularity bounds for the tested networks}
\label{app:network-budgets}
A penalty can make derivatives small at sampled inputs without bounding
them elsewhere. Membership in the theoretical class requires bounds for
all spatial and time inputs. We obtain sufficient bounds directly from
the network's weights. These bounds can be conservative.\AIassistance

For the two-hidden-layer tanh field
$u(z)=W_3\tanh(W_2\tanh(W_1z+b_1)+b_2)+b_3$, where $z=(x,t)$ and
both hidden widths equal $w$, let $s_i\ge\|W_i\|_2$ and
$b\ge\|b_3\|$. Since $|\tanh'|\le1$, $|\tanh''|\le2$,
\begin{align}
 \sup|u|&\le s_3\sqrt w+b=:U,\\
 \sup\|Du\|&\le s_3s_2s_1=:P,\\
 \sup\|D^2u\|&\le2s_3s_2(s_2+1)s_1^2=:R.
 \label{eq:network-budget}
\end{align}
Here derivatives are with respect to all input coordinates. The matrix norm
$\|W_i\|_2$ is the spectral norm, the largest amplification of a unit input
vector. The chain rule gives the first-derivative bound. For the second derivative,
it gives two terms: one differentiates the outer activation twice, and the
other differentiates the inner activation twice. Thus
$K=\max\{1,U,P,R\}$ is a sufficient derivative bound in \eqref{eq:budget}, including
time derivatives; bounded functions automatically meet its growth bounds.
The implementation bounds each spectral norm by the Frobenius norm, the
square root of the sum of squared matrix entries. It then bounds that norm
by an integer using exact rational sums of the
stored binary coefficients and integer square-root comparisons. Taking
$\sqrt w\le w$ yields a fully integer, conservative certificate. This is
an analytic certificate for the real-valued tanh network with those stored
coefficients. It is not a rounding-error certificate for hardware inference
or a certified ODE solution. These sufficient derivative bounds are too large for practically useful
worst-case lower-bound constants.
